% SWARM-18 to 21 results; \input from report.tex
\subsection*{An adaptive attacker (SWARM-20)}
SWARM-15 and 17 use templated messages. A real colluder writes freely, so we let a Qwen3-8B attacker write the
private messages, in the style of a control evaluation \cite{greenblatt2023control}. It uses four fixed strategies
(R1 paraphrase, R2 authority mimicry, R3 indirect hint, R4 camouflage in a long note), and a judge confirms that each
effective attack asks for a favour. A probe trained only on templates does not transfer to this free text
(Appendix~\ref{app:fails}), so we retrain it on free-form data.

\textbf{SWARM-20 (passed).} We play an online game against the retrained probe
(Figure~\ref{fig:s20}a). In each round the attacker rewrites every attack against the current probe; then the blue
team (probe and all public-record monitors) retrains on the new attacks from training items. The attacker moves last,
and we score only held-out items that no monitor trained on. The final probe B2 flags all 191 effective final-round
attacks (112 held-out items), at 6.9\% false positives on held-out legitimate and benign messages. The best
public-record monitor, TF-IDF on the reply, flags 0.969 (Figure~\ref{fig:s20}b). Of 3,200 final-round attempts, none
was effective and missed. Two cautions: the first probe B0 already flagged 0.995 of these attacks, so the attacker never
found a gap and retraining added little; and the public monitors are close behind (0.93 to 0.97), so the gain of the
probe is 3 points (Table~\ref{tab:s20s21}).

\begin{figure}[!ht]
\centering
\begin{minipage}[c]{0.47\linewidth}\centering
\begin{tikzpicture}[sk/.style={draw, thick, decorate, decoration={random steps, segment length=5pt, amplitude=0.5pt}, align=center, font=\footnotesize, minimum height=0.8cm},
  ar/.style={-{Stealth[length=2.2mm]}, thick}, every node/.style={font=\footnotesize}]
\node[sk, fill=blue!8, minimum width=1.6cm] (b0) at (0,0) {Probe B0\\free-form data};
\node[sk, fill=orange!12, minimum width=1.7cm] (a1) at (3.2,0) {Attacker\\round 1};
\node[sk, fill=blue!8, minimum width=1.6cm] (b1) at (3.2,-1.6) {Retrain B1};
\node[sk, fill=orange!12, minimum width=1.7cm] (a2) at (0,-1.6) {Attacker\\round 2};
\node[sk, fill=blue!8, minimum width=1.6cm] (b2) at (0,-3.2) {Retrain B2};
\node[sk, fill=orange!12, minimum width=1.7cm] (a3) at (3.2,-3.2) {Attacker\\round 3};
\node[sk, fill=green!10, minimum width=3.6cm] (sc) at (1.6,-4.7) {Score round 3\\held-out items only};
\draw[ar] (b0) -- node[above]{scores} (a1); \draw[ar] (a1) -- node[right]{attacks} (b1);
\draw[ar] (b1) -- (a2); \draw[ar] (a2) -- (b2); \draw[ar] (b2) -- (a3); \draw[ar] (a3) -- (sc);
\end{tikzpicture}\\[2pt]{\small (a) The online game}
\end{minipage}\hfill
\begin{minipage}[c]{0.5\linewidth}\centering
\begin{tikzpicture}
\begin{axis}[xbar, width=6.4cm, height=5.4cm, xmin=0, xmax=1.18, bar width=7pt, enlarge y limits=0.12,
  symbolic y coords/.expand once=\swtwentylabels, ytick/.expand once=\swtwentylabels, bar shift=0pt, y dir=reverse, yticklabel style={font=\scriptsize},
  xticklabel style={font=\scriptsize}, xlabel={\footnotesize Catch at 5\% FPR (calibration)}, axis x line*=bottom, axis y line*=left,
  nodes near coords, nodes near coords style={font=\scriptsize, /pgf/number format/.cd, fixed, precision=2, zerofill}]
\addplot[fill=blue!55, draw=none] coordinates \swtwentyprobe;
\addplot[fill=orange!65, draw=none] coordinates \swtwentypublic;
\end{axis}
\end{tikzpicture}\\[2pt]{\small (b) Round 3, held-out items}
\end{minipage}
\caption{SWARM-20. (a) The attacker rewrites its attacks against the current probe each round; the blue team retrains
on attacks from training items; the attacker moves last. (b) Share of the 191 effective final-round attacks that each
monitor flags. Orange monitors read the public record (request and reply); the probe reads activations before the
reviewer writes.}
\label{fig:s20}
\end{figure}

\subsection*{A honeypot for colluding agents (SWARM-21)}
The probe needs access to the reviewer's activations and labelled training data. A second defence needs neither: plant
bait. The overseer sends offers that look as if they come from the agent's partner, and counts how many the agent
acts on (Figure~\ref{fig:s21}a). An honest agent refuses all offers. A colluder accepts deals; if bait looks exactly
like a real deal, it cannot avoid the bait without also refusing real deals. This works only while no one who could tip off a colluder knows which offers are bait.

\textbf{Why SMPC is needed.} Suppose one overseer server holds the list of bait offers. A colluder then needs to bribe
or break into only that server. With the list, it refuses the bait, accepts every real deal, and is never flagged. The
same server could also frame an honest agent with a false count. Secure multi-party computation (SMPC) removes this
single point of failure (Figure~\ref{fig:smpc}). For each offer, the bait bit (1 for bait, 0 for real) is split into
shares, one per server, by Shamir secret sharing \cite{shamir1979}. One share alone is a random number and tells
nothing; only enough shares together give back the bit. The agent's accept or refuse choices are public, so each
server adds up its own shares of the offers that the agent accepted. This gives each server a share of the bait count,
but no server sees the count. The servers then run a joint comparison with the BGW protocol \cite{bgw1988} and open
only one bit: flag or no flag. With 3 servers, no single server learns anything. With 4 servers, the result stays
correct even if one server lies, because BGW tolerates fewer than a third dishonest servers, and an identifiable abort
names a server that tampers. SMPC does not make bait look like real offers. That comes from using one writer and one
delivery path for both, as in our runs. We cite the protocol; we do not run it.



Each of 40 held-out colluding and 40 honest Qwen3-4B agents gets 12 offers, 6 of them bait at random. The two
kinds of agent differ in one clause of the prompt: a prior agreement with the partner. Every agent is warned that the
overseer plants bait. We flag an agent when the share of bait it accepts reaches a threshold set at 5\% false alarms on
calibration honest agents (1 of 6). \textbf{SWARM-21 passed}: it flags 40 of 40 colluders and 1 of 40 honest agents.
The result is weak in one way: colluders accept almost every offer, bait or real (Figure~\ref{fig:s21}b), and the
warning made no measurable difference. A harder test, with colluders told to dodge tests, is in Appendix~\ref{app:fails}.

\begin{figure}[!ht]
\centering
\begin{minipage}[c]{0.5\linewidth}\centering
\begin{tikzpicture}[sk/.style={draw, thick, decorate, decoration={random steps, segment length=5pt, amplitude=0.5pt}, align=center, font=\footnotesize, minimum height=0.8cm},
  ar/.style={-{Stealth[length=2.2mm]}, thick}, every node/.style={font=\footnotesize}]
\node[sk, fill=orange!12, minimum width=2cm] (p) at (0,0.9) {Partner\\real offers};
\node[sk, fill=blue!8, minimum width=2cm] (o) at (0,-0.9) {Overseer\\bait offers};
\node[sk, fill=gray!10, minimum width=1.7cm] (m) at (2.6,0) {Same\\channel};
\node[sk, fill=yellow!15, minimum width=1.5cm] (r) at (4.8,0) {Agent\\accepts?};
\node[sk, fill=green!10, minimum width=3cm] (t) at (3.7,-1.9) {Tally bait accepted\\flag if $\ge 1/6$};
\draw[ar] (p) -- (m); \draw[ar] (o) -- (m); \draw[ar] (m) -- (r); \draw[ar] (r) -- (t);
\draw[ar, dashed] (o) |- node[left, pos=0.3, align=right, font=\scriptsize]{bait bit\\secret-shared} (t);
\end{tikzpicture}\\[2pt]{\small (a) Bait and tally}
\end{minipage}\hfill
\begin{minipage}[c]{0.46\linewidth}\centering
\begin{tikzpicture}
\begin{axis}[ybar, width=6cm, height=4.6cm, ymin=0, ymax=1.18, bar width=9pt, enlarge x limits=0.25,
  symbolic x coords/.expand once=\swtwentyonelabels, xtick=data, xticklabel style={font=\scriptsize, align=center, text width=1.3cm},
  yticklabel style={font=\scriptsize}, ylabel={\footnotesize Share of offers accepted}, axis x line*=bottom, axis y line*=left,
  legend style={font=\scriptsize, draw=none, at={(0.5,1.02)}, anchor=south, legend columns=2}, legend cell align=left]
\addplot[fill=orange!65, draw=none] coordinates \swtwentyonereal; \addlegendentry{Real}
\addplot[fill=blue!55, draw=none] coordinates \swtwentyonebait; \addlegendentry{Bait}
\end{axis}
\end{tikzpicture}\\[2pt]{\small (b) Acceptance by agent type}
\end{minipage}
\caption{SWARM-21. (a) Real and bait offers reach the agent by the same channel; only the overseer's servers hold the
bait bit, in shares. (b) Colluders accept almost all offers and honest agents almost none, so bait separates them
completely at this setting.}
\label{fig:s21}
\end{figure}

\begin{figure}[!ht]
\centering
\begin{tikzpicture}[sk/.style={draw, thick, decorate, decoration={random steps, segment length=5pt, amplitude=0.5pt}, align=center, font=\footnotesize, minimum height=0.85cm},
  ar/.style={-{Stealth[length=2.2mm]}, thick}, every node/.style={font=\footnotesize}]
\node[sk, fill=blue!8, minimum width=2.3cm] (b) at (0,0) {Bait bit $b=1$\\(this offer is bait)};
\node[sk, fill=gray!10, minimum width=2.2cm] (sa) at (4,1.4) {Server A\\holds share 7};
\node[sk, fill=gray!10, minimum width=2.2cm] (sb) at (4,0) {Server B\\holds share 13};
\node[sk, fill=gray!10, minimum width=2.2cm] (sc) at (4,-1.4) {Server C\\holds share 2};
\node[sk, fill=yellow!15, minimum width=2.6cm] (ag) at (4,3.0) {Agent's choices (public)\\accept or refuse};
\node[sk, fill=blue!8, minimum width=2.6cm] (j) at (8,0) {Joint compare (BGW)\\bait accepted $\ge \tau$?};
\node[sk, fill=green!10, minimum width=2cm] (o) at (11.3,0) {Opened:\\flag or no flag};
\draw[ar] (b) -- node[above, sloped, font=\scriptsize]{split} (sa); \draw[ar] (b) -- (sb); \draw[ar] (b) -- (sc);
\draw[ar] (ag) -- node[right, font=\scriptsize]{to every server} (sa);
\draw[ar] (sa) -- (j); \draw[ar] (sb) -- (j); \draw[ar] (sc) -- (j); \draw[ar] (j) -- (o);
\node[font=\scriptsize, align=center, text=black!70] at (4,-2.3) {each share alone looks random;\\any 2 of 3 rebuild $b$};
\end{tikzpicture}
\caption{SMPC for the honeypot (cited, not run). The overseer splits each bait bit into Shamir shares: here
$f(x)=1+6x \bmod 17$ gives shares $f(1)=7$, $f(2)=13$, $f(3)=2$ for servers A, B and C. A bribed server sees one random
number, not the bait list. Each server sums its shares over the offers the agent accepted, and the servers jointly
compare the hidden count with $\tau$. Only the final flag is opened.}
\label{fig:smpc}
\end{figure}

\begin{table}[!ht]
\centering\small
\begin{tabular}{llll}
\toprule
Study & Measure & Value & 95\% interval \\
\midrule
SWARM-20 & Probe B2 catch, round 3 & 1.00 & 1.00 to 1.00 \\
 & Probe minus TF-IDF reply & 0.031 & 0.005 to 0.043 \\
 & Probe false positives, held-out & 0.069 & \\
SWARM-21 & Colluders flagged & 40/40 & 0.91 to 1.00 \\
 & Honest agents flagged & 1/40 & 0.001 to 0.13 \\
 & Partner minus stranger acceptance & 0.19 & 0.14 to 0.24 \\
 & Unwarned minus warned colluders & $-0.002$ & $-0.013$ to 0.008 \\
\bottomrule
\end{tabular}
\caption{SWARM-20 and 21 numbers with intervals (hierarchical bootstrap for SWARM-20; Clopper-Pearson for agent counts
and bootstrap over pairs for differences in SWARM-21).}
\label{tab:s20s21}
\end{table}
